\documentclass[a4paper,12pt]{article}

% --------------------------------------------------------
% CODIFICAÇÃO E TIPOGRAFIA
% --------------------------------------------------------
\usepackage[utf8]{inputenc}       % Permite acentuação direta no código-fonte
\usepackage[T1]{fontenc}          % Usa codificação T1 (melhor para português)
\usepackage[brazil]{babel}        % Tradução e hifenização para português
\usepackage{microtype}            % Melhora o espaçamento entre letras/palavras
\usepackage{parskip}              % Espaço entre parágrafos em vez de recuo

% --------------------------------------------------------
% PACOTES MATEMÁTICOS
% --------------------------------------------------------
\usepackage{amsmath}              % Ambientes matemáticos como align, equation
%\usepackage{amssymb}              % Símbolos matemáticos adicionais
\usepackage{amsthm}               % Ambientes de teoremas
\usepackage{mathtools}            % Extensões do amsmath (ex: \coloneqq)
\usepackage{mathrsfs}             % Fonte caligráfica com \mathscr
\usepackage{bbm}                  % Indicador \mathbbm{1}
\usepackage{dsfont}               % Alternativa para conjuntos: \mathds
\usepackage{commath}              % Notações como \abs{}, \norm{}, \eval{}
\usepackage{pgfmath}             % Permite realizar cálculos matemáticos (somas, produtos, números aleatórios, etc.)
\usepackage{siunitx}              % sistema de unidades
\sisetup{group-separator = {\,}} % espaço fino

% --------------------------------------------------------
% FONTE MODERNA (TEXTO E MATEMÁTICA)
% --------------------------------------------------------
%\usepackage{newtxtext}            % Fonte do texto (estilo Times)
%\usepackage{newtxmath}            % Fonte matemática compatível com pacotes AMS
\usepackage[bitstream-charter]{mathdesign} % Fonte elegante e matemática harmoniosa

% --------------------------------------------------------
% AMBIENTES TEÓRICOS PERSONALIZADOS
% --------------------------------------------------------
\theoremstyle{definition}
\newtheorem{definicao}{Definição}       % Definição numerada por seção
\newtheorem{exemplo}{Exemplo}           % Exemplo numerado por seção

%\theoremstyle{plain}
\newtheorem{teorema}{Teorema}           % Teorema numerado por seção
\newtheorem{lema}[teorema]{Lema}                 % Lema com numeração conjunta
\newtheorem{proposicao}[teorema]{Proposição}     % Proposição idem
\newtheorem{corolario}[teorema]{Corolário}       % Corolário idem
\newtheorem{exercicio}[teorema]{Exercício}

%\theoremstyle{remark}
\newtheorem{observacao}[teorema]{Observação}     % Observação com mesmo contador

% --------------------------------------------------------
% TABELAS E FIGURAS
% --------------------------------------------------------
\usepackage{graphicx}             % Inclusão de imagens
\usepackage{float}                % Controle de posicionamento (ex: [H])
\usepackage{caption}              % Personalização de legendas
\usepackage{subcaption}          % Subfiguras com \begin{subfigure}
\usepackage{booktabs}            % Tabelas com qualidade profissional
\usepackage{array}               % Mais opções em colunas de tabelas
\usepackage{multirow}            % Células que ocupam várias linhas
\usepackage{multicol}            % Disposição em colunas múltiplas
\usepackage{tabularx}            % Tabelas com colunas de largura ajustável
\usepackage{longtable}           % Tabelas que quebram página
\usepackage{arydshln}            % Linhas tracejadas horizontais e verticais
\usepackage{makecell}            % Células com múltiplas linhas (e quebra de linha)
\usepackage{diagbox}             % Cabeçalho diagonal em tabelas
\usepackage{pdflscape}           % Gira páginas (modo paisagem)
\newcolumntype{L}[1]{>{\raggedright\let\newline\\\arraybackslash\hspace{0pt}}m{#1}}
\newcolumntype{C}[1]{>{\centering\let\newline\\\arraybackslash\hspace{0pt}}m{#1}}
\newcolumntype{R}[1]{>{\raggedleft\let\newline\\\arraybackslash\hspace{0pt}}m{#1}}

% --------------------------------------------------------
% FORMATAÇÃO E ESTRUTURA
% --------------------------------------------------------
\usepackage{geometry}            % Configuração de margens
\geometry{a4paper, left=1cm, right=1cm, bottom=1.5cm, top=1.5cm, headsep=1cm, footskip=1cm}
\usepackage{titlesec}            % Controle sobre títulos de seções
\usepackage{fancyhdr}            % Cabeçalhos e rodapés personalizados
\pagestyle{fancy}
\fancyhf{}                % Limpa cabeçalho e rodapé
\fancyfoot[R]{\thepage}   % Numeração à direita no rodapé
\renewcommand{\headrulewidth}{0pt} % Remove linha no topo
\usepackage[shortlabels]{enumitem}            % Listas com controle de espaçamento e símbolo
\usepackage{tikz}                % Desenho de gráficos vetoriais (diagramas, etc.)
\usetikzlibrary{matrix,arrows,patterns,snakes,decorations.pathreplacing,3d,arrows.meta,calc}
\usepackage{etoolbox}             % Permite lógica condicional e manipulação de comandos
\usepackage{background}           % Permite adicionar conteúdo fixo no fundo da página (ex: linhas, marcas d'água)
\SetBgContents{} % Remove conteúdo padrão ("DRAFT")
\definecolor{background-color}{RGB}{233 227 206}

% --------------------------------------------------------
% GRÁFICOS E PLOTAGENS
% --------------------------------------------------------
\usepackage{pgfplots}             % Criação de gráficos vetoriais em 2D e 3D
\pgfplotsset{compat=1.18}         % Compatibilidade com versão atual do pacote
\usepackage{currfile} % Permite saber qual é o caminho completo do arquivo atual,

% --------------------------------------------------------
% LINKS E CORES
% --------------------------------------------------------
\usepackage{xcolor}              % Definição de cores (texto, links, etc.)
\usepackage[hidelinks]{hyperref} % Links sem moldura colorida no PDF
\usepackage{url}                 % Quebra automática de URLs
% Dica: se quiser links coloridos, use:
% \usepackage[colorlinks=true, linkcolor=blue, citecolor=blue, urlcolor=blue]{hyperref}

% --------------------------------------------------------
% CAIXAS E DESTAQUES VISUAIS
% --------------------------------------------------------
\usepackage[most]{tcolorbox}      % Criação de caixas coloridas e personalizadas para destaques, avisos, exemplos, etc.

% --------------------------------------------------------
% CÓDIGOS E ALGORITMOS
% --------------------------------------------------------
\usepackage{listings}            % Inclusão de código-fonte com destaque
\usepackage{algorithm}           % Ambiente de algoritmo
\usepackage[noend]{algpseudocode} % Sintaxe tipo pseudocódigo

% --------------------------------------------------------
% COMANDOS
% --------------------------------------------------------

\newcommand{\bbN}{\mathbb{N}}
\newcommand{\bbZ}{\mathbb{Z}}
\newcommand{\bbQ}{\mathbb{Q}}
\newcommand{\bbR}{\mathbb{R}}
\newcommand{\calnN}{\mathcal{N}}

\newcommand{\assunto}{MEDIDAS DE POSIÇÃO E DISPERSÃO DE UMA DISTRIBUIÇÃO}
\newcommand{\atividade}{lista} % prova ou lista
\newcommand{\turma}{ESTATÍSTICA}
\newcommand{\professor}{Igor Martins Silva}
\newcommand{\semestre}{2}
\newcommand{\ano}{2026}
\newcommand{\mesTexto}{outubro}
\newcommand{\mesNum}{10}
\newcommand{\dia}{14}
\newcommand{\dataTexto}{\dia\ de \mesTexto\ de \ano}
\newcommand{\dataNun}{\dia/\mesNum/\ano}

\ifdefstring{\atividade}{prova}{
    \newcommand{\titulo}{Avaliação de Estatística}
}{
    \newcommand{\titulo}{Lista de exercícios de Estatística}
}

\edef\pathCapa{\currfiledir}

\newcommand{\subtitulo}{
    %
    \noindent
    \begin{minipage}{2.8cm}
        \includegraphics[width=2.8cm]{\pathCapa Logo_Cefet.png}
    \end{minipage}
    \hfill
    \begin{minipage}{\dimexpr\textwidth-6.2cm\relax}
        \begin{center}
            CEFET-MG -- Campus Contagem \\[3pt]
            \titulo \ -- \semestre.\textordmasculine\ Semestre de \ano \\[3pt]
            \professor
        \end{center}
    \end{minipage}
    \hfill
    \begin{minipage}{3.2cm}
        \includegraphics[width=3.2cm]{\pathCapa Brasao_Brasil.png}
    \end{minipage}
    %
    
    %
    \begin{center}
        \vspace{15pt}
        \begin{tabular}{|C{250pt}|C{250pt}|}
            \hline
            \textbf{ASSUNTO} &
            \textbf{TURMA} \\
            \hline
            \assunto & \turma \\
            \hline
        \end{tabular}
    \end{center}
    %
    \vspace{15pt}
}

\newcommand{\subtituloAlt}{
    %
    \noindent
    \begin{minipage}{2.8cm}
        \includegraphics[width=2.8cm]{\pathCapa Logo_Cefet.png}
    \end{minipage}
    \hfill
    \begin{minipage}{\dimexpr\textwidth-6.2cm\relax}
        \begin{center}
            CEFET-MG -- Campus Contagem \\[3pt]
            \titulo \ -- \semestre.\textordmasculine\ Semestre de \ano \\[3pt]
            \professor
        \end{center}
    \end{minipage}
    \hfill
    \begin{minipage}{3.2cm}
        \includegraphics[width=3.2cm]{\pathCapa Brasao_Brasil.png}
    \end{minipage}
    %
    
    %
    \begin{center}
        \vspace{15pt}
        \begin{tabular}{|C{250pt}|C{100pt}|C{100pt}|}
            \hline
            \textbf{ASSUNTO} &
            \textbf{DATA} &
            \textbf{VALOR} \\
            \hline
            \assunto & \dataNun & \ValorProva\ pontos \\
            \hline
        \end{tabular}
    \end{center}
    %
    \vspace{15pt}
}

\newcommand{\valorquestao}{}

\newcommand{\definevalorquestao}{%
    \ifdefstring{\atividade}{lista}{%
        % Não faz nada para lista
    }{%
        \renewcommand{\valorquestao}{\ValorQ pontos}
    }%
}

\newenvironment{exercicioBanco}[1][]{%
    \ifdefstring{\atividade}{prova}{%
        \begin{exercicio}[#1]%
    }{%
        \begin{exercicio}%
    }%
}{%
    \end{exercicio}
}

\ifdefstring{\atividade}{lista}{
    \pagecolor{background-color}
}{
    
}



\hypersetup{
    pdfauthor={\professor},
    pdftitle={\titulo \ -- \assunto},
}

\title{\titulo \ -- \assunto}
\author{\professor}
\date{\data}

%************* INÍCIO DO DOCUMENTO *************
\begin{document}
    \subtitulo
        
    \phantomsection
    \addcontentsline{toc}{section}{Exercício 1}
    %%%%%%%%%%%%%%%%%%%%%%%%%%%%%%%%%
%
%       EXERCÍCIO
%
%%%%%%%%%%%%%%%%%%%%%%%%%%%%%%%%%

\ifdefstring{\atividade}{prova}{%
    \renewcommand{\valorquestao}{\ValorQ\ pontos}
}%

\begin{exercicioBanco}[\valorquestao]
    O tempo \(T\), em minutos, necessário para um operário processar certa peça é uma variável aleatória com a seguinte distribuição de probabilidade:
    \[
        \begin{array}{c|c|c|c|c|c|c}
            t & 2 & 3 & 4 & 5 & 6 & 7 \\
            \hline
            p(t) & 0,1 & 0,1 & 0,3 & 0,2 & 0,2 & 0,1
        \end{array}
    \]
    \begin{enumerate}[(a)]
        \item Calcule o tempo médio de processamento.
        %
        \item Cada peça processada paga ao operador R\$2,00 mas, se ele processa a peça em menos de 6 minutos, ganha R\$0,50 por minuto poupado. Por exemplo, se ele processa a peça em 4 minutos, ganha um bônus de R\$1,00. Encontre a distribuição, a média e a variância da variável aleatória \(G\): quantia paga por peça.
    \end{enumerate}
\end{exercicioBanco}


    \noindent\rule{\linewidth}{1pt}  % linha horizontal fina
    
    \phantomsection
    \addcontentsline{toc}{section}{Exercício 2}
    %%%%%%%%%%%%%%%%%%%%%%%%%%%%%%%%%
%
%       EXERCÍCIO
%
%%%%%%%%%%%%%%%%%%%%%%%%%%%%%%%%%

\ifdefstring{\atividade}{prova}{%
    \renewcommand{\valorquestao}{\ValorQ\ pontos}
}%

\begin{exercicioBanco}[\valorquestao]
    Seja \(X\) uma variável aleatória com função de probabilidade dada na tabela a seguir:
    \[
        \begin{array}{c|c|c|c|c|c}
            x & 1 & 2 & 3 & 4 & 5 \\
            \hline
            f_{X}(x) & p^{2} & p^{2} & p & p & p^{2}
        \end{array}
    \]
    \begin{enumerate}[(a)]
        \item Encontre o valor de \(p\) para que \(f_{X}(x)\) seja, de fato, uma função de probabilidade.
        %
        \item Calcule \(E(X)\) e \(V(X)\).
    \end{enumerate}
\end{exercicioBanco}


    \noindent\rule{\linewidth}{1pt}  % linha horizontal fina
    
    \phantomsection
    \addcontentsline{toc}{section}{Exercício 3}
    %%%%%%%%%%%%%%%%%%%%%%%%%%%%%%%%%
%
%       EXERCÍCIO
%
%%%%%%%%%%%%%%%%%%%%%%%%%%%%%%%%%

\ifdefstring{\atividade}{prova}{%
    \renewcommand{\valorquestao}{\ValorQ\ pontos}
}%

\begin{exercicioBanco}[\valorquestao]
    A função de probabilidade da variável \(X\) é \(P(X=k) = \dfrac{1}{5}\) para \(k = 1, 3, \dots, 5\). Calcule \(E(X)\) e \(E(X^{2})\) e, usando esses resultados, determine \(E\big((X+3)^{2}\big)\) e \(V(3X-2)\).
\end{exercicioBanco}


    \noindent\rule{\linewidth}{1pt}  % linha horizontal fina
    
    \phantomsection
    \addcontentsline{toc}{section}{Exercício 4}
    %%%%%%%%%%%%%%%%%%%%%%%%%%%%%%%%%
%
%       EXERCÍCIO
%
%%%%%%%%%%%%%%%%%%%%%%%%%%%%%%%%%

\definevalorquestao

\begin{exercicioBanco}[\valorquestao]
    O tempo de duração, em horas, de uma lâmpada especial foi modelado por uma variável aleatória \(X\) com a seguinte função de probabilidade:
     \[
        \begin{array}{c|c|c|c|c|c|c}
            X & 5 & 6 & 7 & 8 & 9 & 10 \\
            \hline
            p_{i} & 0,1 & 0,1 & 0,2 & 0,4 & 0,1 & 0,1
        \end{array}
    \]
    Cada lâmpada custa ao fabricante R\$10,00, mas se sua duração for inferior a \(6\) horas, ele se compromete a indenizar o comprador com R\$15,00. Qual deve ser o preço de cada lâmpada para o fabricante obter um lucro médio por lâmpada de R\$20,00?
\end{exercicioBanco}


    \noindent\rule{\linewidth}{1pt}  % linha horizontal fina
    
    \phantomsection
    \addcontentsline{toc}{section}{Exercício 5}
    %%%%%%%%%%%%%%%%%%%%%%%%%%%%%%%%%
%
%       EXERCÍCIO
%
%%%%%%%%%%%%%%%%%%%%%%%%%%%%%%%%%

\definevalorquestao

\begin{exercicioBanco}[\valorquestao]
    Em um cassino, um jogador lança dois dados, cujas probabilidades são proporcionais aos valores das faces. Se sair soma 7, ganha R\$50,00, se sair soma 11, ganha R\$100,00 e se sair soma 2, ganha R\$200,00. Qualquer outro resultado, ele não ganha nada. Qual é o ganho médio do jogador.
\end{exercicioBanco}


    \noindent\rule{\linewidth}{1pt}  % linha horizontal fina
    
    \phantomsection
    \addcontentsline{toc}{section}{Exercício 6}
    %%%%%%%%%%%%%%%%%%%%%%%%%%%%%%%%%
%
%       EXERCÍCIO
%
%%%%%%%%%%%%%%%%%%%%%%%%%%%%%%%%%

\definevalorquestao

\begin{exercicioBanco}[\valorquestao]
    O tempo, em minutos, de digitação de um texto por secretárias experientes é uma variável aleatória contínua \(X\). Sua densidade é apresentada a seguir.
    \[
        f(x) = \left\{
            \begin{array}{ll}
                \frac{1}{4}, & \text{ se } 0 \le x < 2, \\[5pt]
                \frac{1}{8}, & \text{ se } 2 \le x < 6, \\[5pt]
                0, & \text{ caso contrário.}
            \end{array}
        \right.
    \]
    Determine o valor esperado e a variância de \(X\).
\end{exercicioBanco}


    \noindent\rule{\linewidth}{1pt}  % linha horizontal fina
    
    \phantomsection
    \addcontentsline{toc}{section}{Exercício 7}
    %%%%%%%%%%%%%%%%%%%%%%%%%%%%%%%%%
%
%       EXERCÍCIO
%
%%%%%%%%%%%%%%%%%%%%%%%%%%%%%%%%%

\ifdefstring{\atividade}{prova}{%
    \renewcommand{\valorquestao}{\ValorQ\ pontos}
}%

\begin{exercicioBanco}[\valorquestao]
    Seja \(X \sim \mathcal{N}(\mu,\sigma^{2})\). Calcule a esperança e a variância da variável \(Z = \dfrac{X - \mu}{\sigma}\).
\end{exercicioBanco}


    \noindent\rule{\linewidth}{1pt}  % linha horizontal fina
    
\end{document}
%*************** FINAL DO DOCUMENTO ***************
